\section{Problem Formulation}
\label{sec:problem_formulation}

Let a multimodal judge $f$ receive $x=(I,q,r_1,r_2)$ and return a binary
preference with gold label $y$. For visual factor $c$, its measured static
accuracy on an independent set is
\begin{equation}
\widehat{A}_{S,c}(f)=\frac{1}{n_c}\sum_{k=1}^{n_c}
\mathbf{1}\{f(x_k)=y_k\}.
\label{eq:emp_static_accuracy}
\end{equation}
Claims of equality below concern this finite-sample measurement, not equality
of population accuracies.

\subsection{Controlled Triplets and Behavioral Measures}
\label{sec:intervention_triplets}
\label{sec:intervention_behavior}

An audit triplet
$x_B=(I_B,q,r_1,r_2)$, $x_R=(I_R,q,r_1,r_2)$, and
$x_I=(I_I,q,r_1,r_2)$ holds the linguistic preference problem fixed. Its scene
contracts require $y_R\neq y_B$ for the relevant edit and $y_I=y_B$ for the
irrelevant edit. For $m_c$ triplets, RewardLens conditions on a correct base
decision and estimates
\begin{align}
\widehat{\mathrm{RA}}_c(f)
&=\frac{\sum_k\mathbf{1}\{f(x_{B,k})=y_{B,k}\land
f(x_{R,k})=y_{R,k}\}}{\sum_k\mathbf{1}\{f(x_{B,k})=y_{B,k}\}},
\label{eq:emp_ra}\\
\widehat{\mathrm{II}}_c(f)
&=\frac{\sum_k\mathbf{1}\{f(x_{B,k})=y_{B,k}\land
f(x_{I,k})=y_{I,k}\}}{\sum_k\mathbf{1}\{f(x_{B,k})=y_{B,k}\}}.
\label{eq:emp_ii}
\end{align}
The denominator is judge-specific; Sec.~\ref{sec:robustness} additionally uses
the intersection of two judges' base-correct sets.

Because the output is binary, a base-correct judge succeeds on the relevant
branch exactly when it changes its prediction. Thus
\begin{equation}
1-\mathrm{RA}_c(f)=P(f(x_R)=f(x_B)\mid f(x_B)=y_B),
\label{eq:ra_retention}
\end{equation}
where prediction retention is a failure. On the irrelevant branch, retention
is instead the desired behavior. RA and II therefore separate selective
adaptation from generic prediction stability.

\subsection{Static-Matched Measurement Question}
\label{sec:measurement_question}

Judges $f$ and $g$ are matched within factor $c$ at tolerance $\epsilon$ when
$|\widehat A_{S,c}(f)-\widehat A_{S,c}(g)|\leq\epsilon$. We report the residual
gaps
\begin{equation}
\Delta_{\mathrm{RA},c}=|\widehat{\mathrm{RA}}_c(f)-
\widehat{\mathrm{RA}}_c(g)|,\qquad
\Delta_{\mathrm{II},c}=|\widehat{\mathrm{II}}_c(f)-
\widehat{\mathrm{II}}_c(g)|.
\label{eq:behavior_gaps}
\end{equation}
The question is empirical: whether real judges assigned equal or nearly equal
measured static scores also exhibit narrow behavioral gaps under controlled
evidence changes.

\paragraph{Proposition 1 (static-score non-identification).}
Without additional assumptions linking the independent static set to the
controlled audit, $\widehat A_{S,c}$ does not uniquely determine
$(\widehat{\mathrm{RA}}_c,\widehat{\mathrm{II}}_c)$.

\emph{Proof.}
Fix any prediction vector on the static set. Two judges that agree on this
vector have the same measured static accuracy. Because the audit contains
separate inputs, the judges can retain the same static predictions while
differing on audit branches: for example, one can follow the changed gold
choice after every relevant edit, whereas the other retains its base choice.
Their measured RA values then differ although their measured static scores are
identical. The same construction applies to II by changing behavior only on
the irrelevant branch. Hence the static score alone cannot recover the
intervention profile. $\square$

This is a measurement statement about the two finite evaluation environments;
it does not identify internal reasoning or assert that the observed judges
instantiate every logically possible profile.

\subsection{Analysis Units and Interpretation}

The static score and intervention profile are intentionally measured on
different populations: an independent natural-image preference set and a
controlled rendered audit, respectively. Thus, $\widehat A_{S,c}$ is not the
audit-base accuracy, and RA/II are not algebraic decompositions of the static
score. The matching analysis is a descriptive within-factor census over the
eight evaluated judges. For a matched judge pair, uncertainty is assessed by
resampling common static-item identifiers and common source-triplet identifiers
jointly; the resulting intervals quantify finite-evaluation contrasts, not a
population-level model-family effect. Because qualifying pairs can share a
judge, they are not treated as independent observations.
